\documentclass[12pt]{article}

\usepackage[portuguese]{babel}
\usepackage[utf8]{inputenc}
\usepackage[T1]{fontenc}
\usepackage{geometry}
\usepackage{array}
\usepackage{tabularx}
\usepackage{colortbl}
\usepackage{xcolor}
\usepackage{fancyhdr}
\usepackage{lastpage}

\geometry{margin=2.5cm}

% Cor cinzenta para cabeçalhos de tabela
\definecolor{tablegray}{gray}{0.9}
\definecolor{primaryblue}{RGB}{0, 51, 102}

% Cabeçalho e rodapé
\pagestyle{fancy}
\fancyhf{}

\fancyhead[L]{Registo de Reunião | EcoRide Solutions}
\fancyhead[R]{\thepage\ de \pageref{LastPage}}

\renewcommand{\headrulewidth}{0.4pt}
\renewcommand{\footrulewidth}{0pt}

\begin{document}

{\noindent\LARGE\bfseries\color{primaryblue} Ata de Reunião | Validação de Requisitos}

\vspace{0.4cm}

\section*{Identificação da Reunião}

\begin{tabularx}{\textwidth}{|l|X|}
\hline
\rowcolor{tablegray}
\textbf{ID Reunião} & \#006 \\ \hline
\textbf{Tipo de reunião} & Validação e clarificação de requisitos \\ \hline
\textbf{Data} & 12 de março de 2026 \\ \hline
\textbf{Hora Início} & 14:00 \\ \hline
\textbf{Hora Fim} & 16:00 \\ \hline
\textbf{Local/Meio} & Sala de Reuniões B \\ \hline
\textbf{Estado} & Finalizada \\ \hline
\end{tabularx}

\vspace{0.8cm}

\section*{Participantes}

\begin{tabularx}{\textwidth}{|l|l|X|}
\hline
\rowcolor{tablegray}
\textbf{Nome} & \textbf{Função} & \textbf{Entidade} \\ \hline
João Martins & Gerente & EcoRide Solutions \\ \hline
Silvina Matagal & Secretária Administrativa & EcoRide Solutions \\ \hline
Ramiro Ramalho & Mecânico & EcoRide Solutions \\ \hline
Luís Soares & Analista & Equipa de Desenvolvimento \\ \hline
Eduardo Fernandes & Analista & Equipa de Desenvolvimento \\ \hline
\end{tabularx}

\vspace{0.8cm}
\hrule % Linha separadora opcional
\vspace{0.2cm}

\section*{Objetivo da Reunião}

A reunião teve como objetivo validar e clarificar os requisitos previamente identificados para o sistema de gestão da oficina de reparação de trotinetes. Pretendeu-se garantir que os requisitos representam corretamente as necessidades da oficina e que estão definidos com um nível de detalhe adequado para permitir a sua implementação pela equipa de desenvolvimento.

Durante a reunião foram analisadas possíveis ambiguidades, problemas de estrutura e lacunas funcionais na documentação existente.

\vspace{0.5cm}

\section*{Principais Problemas Identificados}

Durante a análise dos requisitos verificou-se que alguns apresentam ambiguidades ou falta de detalhe relativamente ao comportamento esperado do sistema. Foram discutidos exemplos relacionados com o registo de danos prévios nas trotinetes, critérios para utilização de peças de baixo valor, definição de taxas de armazenamento e gestão de permissões de utilizadores.

Foi também observado que alguns requisitos são demasiado genéricos, agregando várias funcionalidades distintas, especialmente no que diz respeito à geração de relatórios e análise financeira. Nestes casos foi sugerida a sua decomposição em requisitos mais específicos.

Adicionalmente, identificaram-se alguns requisitos demasiado fragmentados, particularmente na área de gestão de peças e fornecedores, sendo recomendada a sua reorganização.

Foram ainda identificadas funcionalidades relevantes que não estavam formalmente descritas nos requisitos, como comunicação com clientes, gestão de peças defeituosas, integração com sistemas externos de faturação e geração de relatórios técnicos de reparação.

\vspace{0.5cm}

\section*{Decisões Tomadas}

\begin{itemize}
\item Reformular os requisitos que apresentam ambiguidades, introduzindo critérios mais concretos.
\item Dividir requisitos demasiado genéricos em requisitos mais específicos.
\item Reorganizar ou consolidar requisitos relacionados com áreas semelhantes.
\item Avaliar a inclusão das novas funcionalidades identificadas durante a análise.
\item Criar um documento consolidado contendo todos os requisitos revistos do sistema.
\end{itemize}

\vspace{0.5cm}

\section*{Próximos Passos}

\begin{itemize}
\item Revisão e reformulação dos requisitos ambíguos identificados.
\item Reorganização da estrutura geral dos requisitos.
\item Avaliação e eventual integração dos novos requisitos propostos.
\item Elaboração de um documento consolidado de requisitos.
\item Utilização desse documento como base para as fases seguintes do desenvolvimento do sistema.
\end{itemize}

\vspace{0.5cm}
\hrule % Linha separadora opcional
\vspace{0.2cm}

\section*{Encerramento}

Não havendo mais assuntos a tratar, a reunião foi encerrada após a discussão dos pontos referidos, ficando acordado que a equipa de desenvolvimento procederá à revisão da documentação de requisitos com base nas conclusões alcançadas durante a reunião.

\end{document}